\item La población de átomos en un estado de energía $E_n$ en equilibrio térmico a temperatura $T$ se rige por la distribución de Boltzmann:
$$ N = N_g e^{-(E_n - E_g)/k_{\text{B}}T} $$
\begin{enumerate}
    \item Antes de encender el láser, los átomos de neón están en equilibrio térmico a $27.0^\circ\text{C}$. Calcule la razón de equilibrio $N(E_4^*)/N(E_3)$ para las poblaciones involucradas en la figura TP41.1.
    \item Para que el láser funcione, se produce artificialmente una inversión de población. Suponga que se logra una población en el estado superior $E_4^*$ un $2.00\%$ mayor que en el estado $E_3$. Determine a qué temperatura absoluta la distribución de Boltzmann describiría teóricamente esta inversión de población del $+2.00\%$.
    \item Explique por qué situaciones de inversión de población no ocurren naturalmente en equilibrio térmico.
\end{enumerate}